\documentclass{scrreprt}
\usepackage{fhmsegu}
\usepackage{longtable}
\usepackage{booktabs}
\usepackage[all, defaultlines=5]{nowidow}
\usepackage{import}
\usepackage{setspace}

\addbibresource{quellen_methodik.bib}

\DeclareSIUnit{\euro}{\text{\euro}}
\DeclareSIUnit{\year}{a}
\DeclareSIUnit{\kwh}{kWh}
\DeclareSIUnit{\mwh}{MWh}

\makenomenclature
\graphicspath{{images/}}

\setlength{\parindent}{1.5em}
\setlength{\parskip}{1.0em}

\pgfplotsset{compat=1.18}
\pagestyle{scrheadings}
\clearpairofpagestyles

\ihead{\headmark}
\ohead{\getfhmsauthor}
\cfoot{\pagemark}
\pagestyle{scrheadings}
\setheadsepline{0.4pt}
\renewcommand*{\chapterpagestyle}{scrheadings}

\begin{document}

%Titeldefinition
\fhmstitle{Energetische Bewertung virtueller Gebäude anhand von Jahresdauerlinien der Lüftungswärmeverluste nach DIN V 18599}
\fhmssubtitle{Hausarbeit im Fach Energiepotentiale}
\fhmsauthor{Jan Kruse}
\fhmsdate{\today}
\fhmsemail{jan.kruse@kampmann.de}

\fhmsmaketitlepage
\pagestyle{empty}
\clearpage

\chapter*{Abstract}
% TODO: Kurzfassung DE + EN nach Fertigstellung

\pagestyle{scrheadings}
\pagenumbering{gobble}
\begin{spacing}{1.1}
\tableofcontents
\end{spacing}
\clearpage

\pagenumbering{Roman}
\setcounter{page}{0}
\listoffigures
\addcontentsline{toc}{chapter}{Abbildungsverzeichnis}
\listoftables
\addcontentsline{toc}{chapter}{Tabellenverzeichnis}

\par\vspace{1.0em}
\phantomsection
\addcontentsline{toc}{chapter}{Nomenklatur}
\markboth{Nomenklatur}{Nomenklatur}
{\parindent0pt
 \usekomafont{chapter}\usekomafont{chapterprefix}%
 \textsf{\textbf{Nomenklatur}}\par
}
\vspace{0.6em}
{
\renewcommand{\arraystretch}{1.4}
\begin{tabular}{ll}
    $\varPhi_{\mathrm{V}}$ & Lüftungswärmeverlust in \si{\kilo\watt} \\
    $\dot q_{\mathrm{v}}$ & Außenluftvolumenstrom in \si{\cubic\meter\per\hour} \\
    $\theta_{\mathrm{i}}$ & Raum-Solltemperatur in \si{\celsius} \\
    $\theta_{\mathrm{e}}$ & Außenlufttemperatur in \si{\celsius} \\
    $\theta_{\mathrm{v,mech}}$ & Zulufttemperatur-Sollwert der RLT-Anlage in \si{\celsius} \\
    $\rho$ & Dichte der Luft in \si{\kilo\gram\per\cubic\meter} \\
    $c_{\mathrm{p}}$ & spez. Wärmekapazität der Luft in \si{\kilo\joule\per\kilo\gram\per\kelvin} \\
    $Q_{\mathrm{V}}$ & Jahresenergiemenge Lüftung in \si{\kwh} \\
    $q_{\mathrm{h,18°C,12h,mth}}$ & normierter monatlicher Heizenergiekennwert nach DIN V 18599-3 \\
    $g_{\mathrm{h,u/o,mth}}$ & Gradient des monatlichen Heizenergiekennwerts nach DIN V 18599-3 \\
    $f_{\mathrm{T,h}}$ & Korrekturfaktor Betriebszeit (Heizfall) nach DIN V 18599-3 \\
    $\theta_{\mathrm{hc}}$ & Zulufttemperatur-Sollwert im Kennwertverfahren nach DIN V 18599-3 in \si{\celsius} \\
    $t_{\mathrm{v,mech}}$ & tägliche Betriebszeit der RLT-Anlage nach DIN V 18599-3 in \si{\hour} \\
    $d_{\mathrm{v,mech}}$ & jährliche Betriebstage der RLT-Anlage nach DIN V 18599-3 in \si{\day\per\year} \\
    $\eta_{\mathrm{WRG}}$ & Temperaturänderungsgrad der Wärmerückgewinnung (-) \\
    $\theta_{\mathrm{i,c}}$ & Raum-Solltemperatur Kühlung in \si{\celsius} \\
    $q_{\mathrm{c,18°C,12h,mth}}$ & normierter monatlicher Kühlenergiekennwert nach DIN V 18599-3 \\
    $g_{\mathrm{c,u/o,mth}}$ & Gradient des monatlichen Kühlenergiekennwerts nach DIN V 18599-3 \\
    $f_{\mathrm{T,c}}$ & Korrekturfaktor Betriebszeit (Kühlfall) nach DIN V 18599-3 \\
    $x$ & Wasserdampfgehalt der Luft in \si{\gram\per\kilo\gram} \\
    $x_{\mathrm{soll}}$ & Mindest-Wasserdampfgehalt der Zuluft in \si{\gram\per\kilo\gram} \\
    $h$ & spezifische Enthalpie feuchter Luft in \si{\kilo\joule\per\kilo\gram} \\
    $P_{\mathrm{V}}$ & elektrische Ventilatorleistung in \si{\kilo\watt} \\
    $\Delta p_{\mathrm{ZUL/ABL}}$ & Druckerhöhung Zu-/Abluftventilator in \si{\pascal} \\
    $\eta_{\mathrm{V}}$ & Gesamtwirkungsgrad Ventilator/Motor/Antrieb (-) \\
    $f_{\mathrm{p}}$ & Druckverhältnis-Zahl (Teillast) nach DIN V 18599-3 (-) \\
    $W_{\mathrm{V}}$ & Jahresenergiemenge Ventilatoren in \si{\kwh} \\
    % TODO: weitere Symbole ergänzen
\end{tabular}
}

\addchap{Abkürzungsverzeichnis}
\begin{acronym}[ASHRAE]
    \acro{RLT}{Raumlufttechnik}
    \acro{NWG}{Nichtwohngebäude}
    \acro{TRY}{Testreferenzjahr}
    \acro{TRJ}{Testreferenzjahr-Datensatz}
    \acro{DWD}{Deutscher Wetterdienst}
    \acro{BBR}{Bundesamt für Bauwesen und Raumordnung}
    \acro{JDL}{Jahresdauerlinie}
    \acro{WRG}{Wärmerückgewinnung}
    \acro{SFP}{Specific Fan Power}
    \acro{VAV}{Variable Air Volume}
    % TODO: weitere Abkürzungen bei Einbau in Fließtext ergänzen
\end{acronym}

\pagestyle{scrheadings}

\chapter{Einleitung}
\pagenumbering{arabic}
% Motivation: Anteil Lüftungswärmeverluste am Energiebedarf von NWG
% Ziel: Python-Tool zur überschlägigen Jahresenergiemenge Lüftung + JDL
% Aufbau der Arbeit (Verweis auf Kap. 2-5)

\chapter{Theoretischer Hintergrund}\label{chap:Theorie}

\section{Energetische Bewertung von Gebäuden nach DIN V 18599}
Die Vornormenreihe DIN V 18599 regelt die energetische Bewertung von
Gebäuden hinsichtlich des Nutz-, End- und Primärenergiebedarfs für
Heizung, Kühlung, Lüftung, Trinkwarmwasser und Beleuchtung
\parencite{DIN1859910}. Sie gliedert sich in zwölf aufeinander
aufbauende Teile: Teil 1 stellt das allgemeine Bilanzierungsverfahren
und die Verknüpfung der übrigen Teile bereit, Teil 2 behandelt den
Nutzenergiebedarf für Heizen und Kühlen von Gebäudezonen, Teil 3 den
Nutzenergiebedarf für die energetische Luftaufbereitung, Teil 7 den
Endenergiebedarf von \ac{RLT}- und Klimakältesystemen für den
Nichtwohnungsbau sowie Teil 10 die Nutzungsrandbedingungen und
Klimadaten \parencite{DIN1859910}. Die energetische Bilanzierung
erfolgt dabei grundsätzlich monatsweise auf Basis eines
Referenzklimas für Deutschland \parencite{DIN185993}.

\section{Das Kennwertverfahren nach DIN V 18599-3}
Für die Ermittlung des Nutzenergiebedarfs raumlufttechnischer Anlagen
stellt DIN V 18599-3 ein sogenanntes Kennwertverfahren bereit
\parencite{DIN185993}. Diesem liegt folgender Grundgedanke zugrunde:
Für die in der Praxis häufigsten \ac{RLT}-Systeme wurde eine
umfangreiche Variantenmatrix aus 46 Anlagenvarianten gebildet, die
sich hinsichtlich Feuchteanforderung, Befeuchtertyp,
\ac{WRG}-Typ und \ac{WRG}-Größe unterscheiden \parencite{DIN185993}.
Für jede dieser Varianten wurde der Nutzenergiebedarf vorab in einer
stündlichen Simulation für einen festgelegten Basisfall (Referenzklima,
Zulufttemperatur-Sollwert $\theta_{\mathrm{v,mech}}=\SI{18}{\celsius}$,
tägliche Betriebszeit $t_{\mathrm{v,mech}}=\SI{12}{\hour}$) ermittelt und
als normierter, monatlicher Energiekennwert (z.\,B.
$q_{\mathrm{h,18°C,12h,mth}}$) tabellarisch hinterlegt
\parencite{DIN185993}.

Zur Anwendung auf ein konkretes Gebäude wird aus der Variantenmatrix
die zutreffende Anlagenvariante ausgewählt und der zugehörige
Kennwert mithilfe von Gradienten (für abweichende
Zulufttemperatur-Sollwerte) sowie Korrekturfaktoren (für abweichende
tägliche Betriebszeiten) auf die tatsächlichen Randbedingungen
umgerechnet und auf den tatsächlichen Volumenstrom denormiert
\parencite{DIN185993}. Der Wärmerückgewinnungsgrad einer Anlage kann
dabei durch lineare Interpolation zwischen den in der Variantenmatrix
hinterlegten Stützstellen (45\,\%, 60\,\%, 75\,\%) berücksichtigt
werden \parencite{DIN185993}. Das Verfahren erlaubt damit eine
monatsweise, überschlägige Bilanzierung, ohne dass für jedes Gebäude
eine vollständige stündliche Simulation der Anlagentechnik
durchgeführt werden muss. Alternativ lässt die Norm für nicht durch
die Variantenmatrix abgebildete Anlagenkonzepte auch ein direktes
Stundenschrittverfahren zu \parencite{DIN185993}.

\section{Testreferenzjahre als Wetterdatengrundlage}
Für die energetische Bewertung von Gebäuden und ihrer technischen
Anlagen werden repräsentative, stündlich aufgelöste Wetterdatensätze
benötigt, die sogenannten \acp{TRY}. Diese liefern für ein
vollständiges Jahr klimatologische Randbedingungen, ohne die
Extremwerte eines einzelnen realen Messjahres widerzuspiegeln
\parencite{DWDBBR2017}. Während frühere Testreferenzjahre auf 15
Klimaregionen mit jeweils einer Repräsentanzstation beruhten
\parencite[vgl.][]{DIN1859910}, wurden diese im Rahmen der
TRY-Weiterentwicklung durch ortsgenaue Datensätze mit einer
räumlichen Auflösung von etwa \SI{1}{\square\kilo\meter} für ganz
Deutschland abgelöst \parencite{Spekat2011,DWDBBR2017}. Die
Interpolation der Basisgrößen (u.\,a. Lufttemperatur,
Windgeschwindigkeit, kurz- und langwellige Strahlung) auf dieses
Zielraster erfolgt unter Verwendung von Stations-, Satelliten- und
Modelldaten, u.\,a. des Regionalklimamodells COSMO-CLM
\parencite{Krähenmann2016,DWDBBR2017}.

Ein mittleres \ac{TRJ} wird dabei nicht künstlich erzeugt, sondern
aus realen, mehrtägigen Witterungsabschnitten des Zeitraums 1995 bis
2012 zusammengesetzt, die so ausgewählt werden, dass Mittelwert und
Streuung der stündlichen Lufttemperatur sowie der Mittelwert der
Globalstrahlung im Vergleich zum langjährigen Jahresgang
bestmöglich getroffen werden \parencite{DWDBBR2017}. Neben den
mittleren \acp{TRJ} werden zusätzlich extreme \acp{TRJ} (ein
sommer-fokussiertes, extrem warmes und ein winter-fokussiertes,
extrem kaltes Jahr) sowie Zukunfts-\acp{TRJ} auf Basis regionaler
Klimaprojektionen für den Zeitraum 2031 bis 2060 bereitgestellt
\parencite{DWDBBR2017}. Für die vorliegende Arbeit wird ausschließlich
das mittlere, gegenwärtige \ac{TRJ} verwendet, da eine überschlägige
Jahresenergiemenge unter heutigen, repräsentativen
Witterungsbedingungen ermittelt werden soll.

\section{Lüftungswärmeverluste und Jahresdauerlinie}
% Definition Lüftungswärmeverlust (Gleichung), Mindestaußenluftvolumenstrom
% Definition Jahresdauerlinie und Nutzen zur energetischen Bewertung

\chapter{Methodik}\label{chap:Methodik}

Zur überschlägigen Ermittlung der Jahresenergiemenge der Lüftung und zur
Erstellung der zugehörigen \ac{JDL} wurde im Rahmen dieser Arbeit ein
Python-Werkzeug entwickelt. Es berechnet auf Basis stündlicher
Wetterdaten und definierter Nutzungsrandbedingungen den
Lüftungswärmeverlust virtueller \ac{NWG} und stellt diesen als
Jahresdauerlinie dar.

\section{Aufbau des Berechnungstools}
Das Tool ist als Python-Paket \texttt{ventilation\_jdl} mit folgender
Modulstruktur aufgebaut:

\begin{itemize}
    \item \texttt{weather.py}: Einlesen der \ac{TRY}-Datensätze
    \item \texttt{building.py}: Definition virtueller Gebäude und
        Nutzungsrandbedingungen
    \item \texttt{ventilation.py}: stündliche Berechnung des
        Lüftungswärmeverlusts
    \item \texttt{load\_duration.py}: Sortierung und Darstellung der
        \ac{JDL}
    \item \texttt{kennwertverfahren.py}: Vergleichsrechnung nach dem
        Kennwertverfahren der DIN V 18599-3 (Abschnitt~\ref{sec:Vergleichsrechnung})
    \item \texttt{cli.py}: Kommandozeilen-Schnittstelle zur
        Zusammenführung der Berechnungsschritte
\end{itemize}

Als Eingabeparameter dienen eine \ac{TRY}-Datei sowie ein
Gebäudeprofil (Außenluftvolumenstrom, Raum-Solltemperatur,
Betriebszeiten, Wärmerückgewinnungsgrad) in Form einer JSON-Datei.
Ausgabegrößen sind die Jahresenergiemenge Lüftung $Q_{\mathrm{V}}$, die
rechnerischen Volllaststunden sowie die \ac{JDL} als Diagramm und
CSV-Datei. Die Datenverarbeitung erfolgt mit den Python-Bibliotheken
\texttt{pandas} und \texttt{numpy}, die grafische Auswertung mit
\texttt{matplotlib}. Die Berechnungslogik ist durch automatisierte
Tests (\texttt{pytest}) auf Basis synthetischer Wetterdaten abgesichert.

\section{Wetterdaten}
Als Wetterdatenbasis werden die ortsgenauen \acp{TRY} des \ac{DWD}
verwendet, die 2017 im Auftrag des \ac{BBR} in Zusammenarbeit mit dem
\ac{DWD} veröffentlicht wurden \parencite{DWDBBR2017}. Im Unterschied
zu den älteren, auf 15 Klimaregionen mit Repräsentanzstationen
basierenden Testreferenzjahren nach DIN V 18599-10, Anhang E
\parencite[vgl.][]{DIN1859910}, liegen die ortsgenauen \acp{TRY} für
ein Raster von etwa \SI{1}{\square\kilo\meter} über ganz Deutschland
vor \parencite{DWDBBR2017}. Jeder Datensatz enthält 8760 Stundenwerte
meteorologischer Größen (u.\,a. Lufttemperatur, Windgeschwindigkeit,
Global- und Wärmestrahlung) für ein Jahr \parencite{DWDBBR2017}. Für
das Berechnungstool wird ausschließlich die stündliche
Außenlufttemperatur $\theta_{\mathrm{e}}$ ausgewertet.

Die TRY-Rohdaten liegen als ASCII-Dateien mit einem Metadaten-Header,
einer Zeile mit Parameterabkürzungen als Spaltennamen und 8760
Datenzeilen vor. Da sich die genaue Spaltenreihenfolge zwischen
TRY-Varianten leicht unterscheiden kann, sucht das Modul
\texttt{weather.py} die Spaltennamenzeile in der Datei selbst
(anhand bekannter Parameterabkürzungen wie \texttt{RW}, \texttt{HW},
\texttt{MM}, \texttt{DD}, \texttt{HH}, \texttt{t}), anstatt eine feste
Spaltenreihenfolge anzunehmen. Dadurch bleibt das Tool robust
gegenüber unterschiedlichen TRY-Dateiformaten.

\section{Definition der virtuellen Gebäude}
Die virtuellen Gebäude werden anhand der Nutzungsrandbedingungen für
\ac{NWG} nach DIN V 18599-10 definiert \parencite{DIN1859910}. Für
jedes virtuelle Gebäude werden das zugeordnete Nutzungsprofil (z.\,B.
Einzelbüro, Lagerhalle), der Außenluftvolumenstrom, die
Raum-Solltemperatur, die täglichen Betriebszeiten sowie ein
Wärmerückgewinnungsgrad hinterlegt. Die konkreten Zahlenwerte werden
den Tabellen der DIN V 18599-10 entnommen bzw. für die untersuchten
Varianten projektspezifisch festgelegt. Im Berechnungstool werden
diese Randbedingungen als JSON-Profil abgelegt (siehe
\texttt{data/profiles/}), sodass weitere Gebäudevarianten ohne
Codeänderung ergänzt werden können.

\section{Berechnung der Jahresenergiemenge Lüftung}
Die energetische Bewertung von Gebäuden nach DIN V 18599 erfolgt im
Regelfall über ein monatliches Bilanzverfahren
\parencite{DIN1859910}; für die Lüftungswärme- und -kältebedarfe
raumlufttechnischer Anlagen stellt DIN V 18599-3 hierzu das in
Abschnitt~\ref{chap:Theorie} beschriebene Kennwertverfahren bereit,
das auf einer Variantenmatrix von 46 RLT-Anlagen-Varianten mit
monatlich normierten Energiekennwerten sowie Korrekturfaktoren für
Betriebszeit, Zulufttemperatur und Wärmerückgewinnungsgrad beruht
\parencite{DIN185993}.

Für die vorliegende, überschlägige Betrachtung wird stattdessen eine
vereinfachte stündliche Energiebilanz auf Basis der \ac{TRY}-Daten
angesetzt. Diese folgt dem Grundprinzip der Lüftungswärmeverlustbilanz
nach DIN V 18599-2 \parencite{DIN185992}, wird jedoch nicht mit
Monatsmittelwerten, sondern mit stündlichen Temperaturwerten
ausgewertet, um im Anschluss eine \ac{JDL} bilden zu können. Der
Lüftungswärmeverlust einer Stunde $t$ berechnet sich zu

\begin{equation}\label{eq:QdotV}
    \dot Q_{\mathrm{V}}(t) = \rho \, c_{\mathrm{p}} \, \dot q_{\mathrm{v}}(t) \, \Delta\theta_{\mathrm{eff}}(t)
\end{equation}

mit

\begin{equation}\label{eq:dtheta}
    \Delta\theta_{\mathrm{eff}}(t) =
    \begin{cases}
        (1 - \eta_{\mathrm{WRG}}) \, \bigl(\theta_{\mathrm{i}} - \theta_{\mathrm{e}}(t)\bigr) & \text{für } \theta_{\mathrm{e}}(t) < \theta_{\mathrm{i}} \text{ und Betriebszeit} \\
        0 & \text{sonst.}
    \end{cases}
\end{equation}

Der Außenluftvolumenstrom $\dot q_{\mathrm{v}}(t)$ wird nur innerhalb
der hinterlegten Betriebszeit des jeweiligen virtuellen Gebäudes
angesetzt, eine vorhandene \ac{WRG} wird über den
Temperaturänderungsgrad $\eta_{\mathrm{WRG}}$ vereinfacht durch
Reduktion der wirksamen Temperaturdifferenz berücksichtigt. Die
Jahresenergiemenge Lüftung $Q_{\mathrm{V}}$ ergibt sich als Summe der
8760 Stundenwerte nach Gleichung \eqref{eq:QdotV}.

Diese Vereinfachung unterscheidet sich in zwei wesentlichen Punkten
vom Kennwertverfahren nach DIN V 18599-3: Erstens erfolgt die
Bilanzierung stündlich statt monatlich, wodurch auf eine gesonderte
Umrechnung über Gradienten $g_{\mathrm{h,u/o,mth}}$
\parencite{DIN185993} verzichtet werden kann, da die tatsächliche
Zulufttemperatur stündlich in die Bilanz eingeht. Zweitens wird der
Einfluss der täglichen Betriebszeit nicht über einen empirischen
Korrekturfaktor $f_{\mathrm{T,h}}$ \parencite{DIN185993} abgebildet,
sondern unmittelbar über die stundengenaue Betriebszeitmaske
berücksichtigt.

Analog zum Heizfall wird ein sensibler Kühlfall berechnet, sofern für
das virtuelle Gebäude eine Kühl-Solltemperatur $\theta_{\mathrm{i,c}}$
hinterlegt ist:

\begin{equation}\label{eq:QdotC}
    \dot Q_{\mathrm{C}}(t) = \rho \, c_{\mathrm{p}} \, \dot q_{\mathrm{v}}(t) \, (1-\eta_{\mathrm{WRG}}) \, \bigl(\theta_{\mathrm{e}}(t) - \theta_{\mathrm{i,c}}\bigr)
    \quad \text{für } \theta_{\mathrm{e}}(t) > \theta_{\mathrm{i,c}} \text{ und Betriebszeit.}
\end{equation}

Zusätzlich wird eine latente Befeuchtungsbilanz gebildet, sofern für
das virtuelle Gebäude ein Mindest-Wasserdampfgehalt der Zuluft
$x_{\mathrm{soll}}$ hinterlegt ist:

\begin{equation}\label{eq:Qst}
    \dot Q_{\mathrm{st}}(t) = \dot m_{\mathrm{tL}} \, \bigl(h(\theta_{\mathrm{i}}, x_{\mathrm{soll}}) - h(\theta_{\mathrm{i}}, x_{\mathrm{e}}(t))\bigr)
    \quad \text{für } x_{\mathrm{e}}(t) < x_{\mathrm{soll}} \text{ und Betriebszeit,}
\end{equation}

wobei $h$ die spezifische Enthalpie feuchter Luft je kg trockener Luft
bezeichnet. Der Wasserdampfgehalt der Außenluft $x_{\mathrm{e}}(t)$
wird direkt der \ac{TRY}-Datei entnommen; die Enthalpiewerte werden
mit der Stoffdatenbibliothek CoolProp \parencite{Bell2014} berechnet,
um belastbare Stoffwerte feuchter Luft ohne eigene, fehleranfällige
Psychrometrie-Näherungsformeln zu erhalten.

Neben dem Nutzenergiebedarf für Heizen, Kühlen und Befeuchten wird
zusätzlich der elektrische Endenergiebedarf der Zu- und
Abluftventilatoren abgeschätzt. Nach DIN V 18599-3 berechnet sich die
elektrische Ventilatorleistung allgemein aus dem geförderten
Volumenstrom, der Druckerhöhung, dem Ventilatorwirkungsgrad sowie
einer Druckverhältnis-Zahl $f_{\mathrm{p}}$, die den Einfluss einer
Volumenstromregelung (\ac{VAV}) bei Teillast abbildet
\parencite{DIN185993}. Für Anlagen mit konstantem Volumenstrom, wie
sie den in dieser Arbeit betrachteten virtuellen Gebäuden zugrunde
liegen, entfällt die Teillastregelung ($f_{\mathrm{p}}=1$), sodass
sich die Gleichung auf den bekannten Sonderfall

\begin{equation}\label{eq:PV}
    P_{\mathrm{V}}(t) = \dot q_{\mathrm{v}}(t) \, \frac{\Delta p_{\mathrm{ZUL}} + \Delta p_{\mathrm{ABL}}}{\eta_{\mathrm{V}}}
    \quad \text{für Betriebszeit,}
\end{equation}

reduziert \parencite{DIN185993}. Der Wirkungsgrad $\eta_{\mathrm{V}}$
fasst dabei, analog zum \ac{SFP}-Wert nach DIN V 18599-7,
Abschnitt~3.1.2 (spezifische elektrische Leistungsaufnahme je
Volumenstrom), Ventilator-, Motor- und ggf. Antriebswirkungsgrad
zusammen \parencite{DIN185997}. Die Jahresenergiemenge der
Ventilatoren $W_{\mathrm{V}}$ ergibt sich, analog zu Gleichung~(25)
der DIN V 18599-3, als Summe der stündlichen Leistungswerte nach
Gleichung~\eqref{eq:PV} über die Betriebsstunden des Jahres
\parencite{DIN185993}. Die für die virtuellen Gebäude angesetzten
Druckerhöhungen und Wirkungsgrade sind projektspezifische, plausible
Auslegungswerte und keine Normwerte, da DIN V 18599-7 hierfür primär
Kennwerte für die anlagentechnische Bewertung von Bestandsanlagen
(Anhang~F) bereitstellt \parencite{DIN185997}.

Diese Erweiterungen sind für eine überschlägige Betrachtung
vertretbar, vernachlässigen jedoch weiterhin anlagenspezifische
Regelungseffekte sowie – im Fall der Befeuchtung – die Unterscheidung
zwischen Dampf- und Verdunstungsbefeuchtung, die im vollständigen
Kennwertverfahren über die Variantenmatrix erfasst werden
\parencite{DIN185993}.

\section{Erstellung der Jahresdauerlinie}
Zur Erstellung der \ac{JDL} werden die 8760 Stundenwerte des
Lüftungswärmeverlusts nach Gleichung \eqref{eq:QdotV} absteigend
sortiert. Die Abszisse der \ac{JDL} gibt dabei die Anzahl der
Jahresstunden an, in denen der jeweilige Leistungswert erreicht oder
überschritten wird. Aus der \ac{JDL} lassen sich zusätzlich die
rechnerischen Volllaststunden als Quotient aus Jahresenergiemenge und
Spitzenlast ableiten.

\section{Vergleichsrechnung mit dem Kennwertverfahren nach DIN V 18599-3}\label{sec:Vergleichsrechnung}
Um die Ergebnisse des stündlichen Bilanzverfahrens einordnen zu
können, wird für das Einzelbüro-Gebäudeprofil zusätzlich eine
Vergleichsrechnung nach dem normativen Kennwertverfahren der DIN V
18599-3 (Jahresverfahren für Anlagen mit konstantem Volumenstrom,
Abschnitt~4.4.1 i.\,V.\,m. Tabelle~A.1) durchgeführt
\parencite{DIN185993}. Diese Gegenüberstellung erfolgt exemplarisch
für ein Gebäude und eine Anlagenvariante und erhebt keinen Anspruch
auf eine vollständige Anwendung der gesamten Variantenmatrix.

Aus der Variantenmatrix nach DIN V 18599-3, Tabelle~7, wird zunächst
die zum Einzelbüro-Profil passende Anlagenvariante identifiziert. Das
Profil sieht keine Feuchteanforderung sowie einen
Wärmerückgewinnungsgrad von $\eta_{\mathrm{WRG}}=0{,}6$ vor; dies
entspricht der Variante~3 (keine Feuchteanforderung,
\ac{WRG}-Typ „nur Wärme“, Wärmerückgewinnungsgrad $60\,\%$)
\parencite{DIN185993}. Für diese Variante weist Tabelle~A.1
(Gesamtjahr) einen Heizenergiekennwert von
$q_{\mathrm{h,18°C,12h}}=\SI{1148}{\watthour\per(\cubic\meter\per\hour)}$
sowie die Gradienten $g_{\mathrm{h,u}}=\SI{274}{\watthour\per(\kelvin\cubic\meter\per\hour)}$
(für $\theta_{\mathrm{hc}}<\SI{18}{\celsius}$) und
$g_{\mathrm{h,o}}=\SI{783}{\watthour\per(\kelvin\cubic\meter\per\hour)}$
(für $\theta_{\mathrm{hc}}>\SI{18}{\celsius}$) aus
\parencite{DIN185993}.

Für das Einzelbüro-Profil ($\theta_{\mathrm{hc}}=\SI{20}{\celsius}$,
tägliche Betriebszeit $t_{\mathrm{v,mech}}=\SI{11}{\hour}$
[7--18\,Uhr], $d_{\mathrm{v,mech}}=\SI{260}{\day\per\year}$
[Mo--Fr], $\dot q_{\mathrm{v}}=\SI{500}{\cubic\meter\per\hour}$)
ergibt sich der Heizenergiekennwert nach Umrechnung auf die
abweichende Zulufttemperatur zu

\begin{equation}
    q_{\mathrm{h}}(\theta_{\mathrm{hc}}) = q_{\mathrm{h,18°C,12h}} + g_{\mathrm{h,o}} \, (\theta_{\mathrm{hc}} - \SI{18}{\celsius}) \approx \SI{2714}{\watthour\per(\cubic\meter\per\hour)}
\end{equation}

Die Umrechnung auf die tatsächliche tägliche Betriebszeit erfolgt
über den Korrekturfaktor $f_{\mathrm{T,h}}$, der für
$t_{\mathrm{v,mech}}=\SI{11}{\hour}$ nahe $1{,}0$ liegt
\parencite{DIN185993}. Nach Skalierung auf die tägliche Betriebszeit,
die tatsächliche Anzahl jährlicher Betriebstage (Tabelle~A.1 basiert
auf $365$~Tagen) sowie Denormierung auf den Außenluftvolumenstrom von
$\SI{500}{\cubic\meter\per\hour}$ ergibt sich ein jährlicher
Nutzenergiebedarf Heizen von rund $Q_{\mathrm{h,Kennwert}}\approx
\SI{880}{\kwh\per\year}$. Die Berechnung wurde im Rahmen des
Python-Tools als eigenständiges Modul
(\texttt{kennwertverfahren.py}) nachgebildet und durch einen Test
gegen den von Hand nachvollzogenen Wert abgesichert.

Analog wird für den Kühlfall der Kühlenergiekennwert
$q_{\mathrm{c,18°C,12h}}=\SI{2309}{\watthour\per(\cubic\meter\per\hour)}$
sowie die Gradienten $g_{\mathrm{c,u}}=\SI{856}{\watthour\per(\kelvin\cubic\meter\per\hour)}$
und $g_{\mathrm{c,o}}=\SI{389}{\watthour\per(\kelvin\cubic\meter\per\hour)}$
der Variante~3 herangezogen \parencite{DIN185993}. Die reale
Raum-Kühlsolltemperatur des Einzelbüro-Profils nach DIN V 18599-10
beträgt $\theta_{\mathrm{i,c}}=\SI{24}{\celsius}$
\parencite{DIN1859910}; dieser Wert liegt jedoch außerhalb des in
Gleichung~(31)/(32) der DIN V 18599-3 festgelegten
Gültigkeitsbereichs von $\SI{14}{\celsius}$ bis $\SI{22}{\celsius}$
für den Zulufttemperatur-Sollwert \parencite{DIN185993}. Um eine
unzulässige Extrapolation zu vermeiden, wird für die
Vergleichsrechnung ersatzweise der obere Rand des Gültigkeitsbereichs
($\SI{22}{\celsius}$) angesetzt. Mit dem zugehörigen Korrekturfaktor
$f_{\mathrm{T,c}}$ für $t_{\mathrm{v,mech}}=\SI{11}{\hour}$
\parencite{DIN185993} ergibt sich nach gleichem Rechengang wie beim
Heizfall ein jährlicher Nutzenergiebedarf Kühlen von rund
$Q_{\mathrm{c,Kennwert}}\approx\SI{249}{\kwh\per\year}$. Für die
Befeuchtung ist innerhalb der Variante~3 kein Vergleich möglich, da
diese Variante keine Feuchteanforderung vorsieht und der zugehörige
Dampfkennwert $q_{\mathrm{st,12h}}=0$ ist \parencite{DIN185993}; ein
normkonformer Vergleich der im Tool implementierten latenten
Befeuchtungsbilanz würde eine Anlagenvariante mit Feuchteanforderung
voraussetzen und wird im Rahmen dieser Arbeit nicht weiter verfolgt.

\chapter{Ergebnisse \& Diskussion}\label{chap:Ergebnisse}

\section{Jahresdauerlinien der virtuellen Gebäude}

\section{Sensitivitätsbetrachtung}
% z. B. Einfluss Luftwechsel, Standort, WRG-Grad

\section{Einordnung und Grenzen der überschlägigen Methode}
Der Vergleich des stündlichen Bilanzverfahrens mit dem
Kennwertverfahren nach DIN V 18599-3 (Abschnitt~\ref{sec:Vergleichsrechnung})
zeigt, in welcher Größenordnung sich die überschlägige Methode
gegenüber dem normativen Verfahren bewegt. % TODO: Vergleich mit tatsächlichem TRY-Tool-Ergebnis ergänzen, sobald ein reales TRY für den betrachteten Standort ausgewertet wurde.

Bei der Interpretation der Abweichung ist zu berücksichtigen, dass
die in DIN V 18599-3, Tabelle~A.1 hinterlegten Energiekennwerte
selbst auf Grundlage eines Testreferenzjahres für die
„Region~4“ des früheren, auf 15 Klimaregionen basierenden
Referenzklimasystems nach DIN V 18599-10, Anhang~E ermittelt wurden
\parencite{DIN185993,DIN1859910}. Das im Rahmen dieser Arbeit
verwendete Berechnungstool greift hingegen auf die seit 2017
verfügbaren, ortsgenauen \acp{TRY} zurück
\parencite{DWDBBR2017,Spekat2011}. Eine Abweichung zwischen beiden
Verfahren ist daher nicht ausschließlich auf die methodische
Vereinfachung der stündlichen Bilanz zurückzuführen, sondern auch auf
eine unterschiedliche Wetterdatengrundlage. Für eine methodisch
sauberere Gegenüberstellung müssten beide Verfahren auf denselben
Klimadatensatz angewendet werden.

Darüber hinaus vernachlässigt das im Tool implementierte
Bilanzverfahren, wie in Abschnitt~\ref{chap:Methodik} dargelegt, den
latenten Wärmeanteil sowie anlagenspezifische Regelungseffekte, die
im Kennwertverfahren über die Variantenmatrix und die zugehörigen
Korrekturfaktoren implizit erfasst werden. Die überschlägige Methode
ist damit als Orientierungsrechnung zu verstehen, die für einen
schnellen Vergleich mehrerer Gebäudevarianten und Standorte sowie zur
Erstellung einer \ac{JDL} geeignet ist, jedoch keinen Ersatz für eine
normkonforme Nachweisführung nach DIN V 18599 darstellt.

Für den Kühlfall kommt als zusätzliche Einschränkung hinzu, dass die
reale Raum-Kühlsolltemperatur des betrachteten Gebäudeprofils
($\SI{24}{\celsius}$) außerhalb des Gültigkeitsbereichs der
Umrechnungsgleichungen des Kennwertverfahrens liegt
(Abschnitt~\ref{sec:Vergleichsrechnung}); die Vergleichsrechnung
Kühlen ist daher nur eingeschränkt aussagekräftig. Für die
Befeuchtung wurde im Tool zwar eine latente Bilanz auf Basis
realer Wasserdampfgehalte implementiert, ein normkonformer Abgleich
mit dem Kennwertverfahren konnte im Rahmen dieser Arbeit mangels
einer geeigneten, Feuchteanforderungen berücksichtigenden
Anlagenvariante für das betrachtete Gebäude nicht durchgeführt
werden.

Für die berechnete Jahresenergiemenge der Ventilatoren gilt
einschränkend, dass die verwendeten Druckerhöhungen und der
Ventilatorwirkungsgrad projektspezifisch angenommen und nicht aus
DIN V 18599-7 entnommen wurden, da die Norm hierfür primär auf
anlagenspezifische Bestandsdaten (Inspektionswerte) verweist
\parencite{DIN185997}. Zudem wird mit $f_{\mathrm{p}}=1$ nur der
Sonderfall konstanten Volumenstroms abgebildet; für Anlagen mit
\ac{VAV}-Teillastregelung wäre die vollständige Gleichung nach
DIN V 18599-3 heranzuziehen \parencite{DIN185993}.

\chapter{Fazit \& Ausblick}\label{chap:Fazit}

\printbibliography

\appendix
\chapter{Parameterübersicht virtuelle Gebäude}
\chapter{Programmübersicht Python-Tool}

\end{document}
